\section*{Capítulo \ref{ch:g4:adapted-to-their-world} --- Adaptados ao mundo deles}

\begin{solution}{exo:g4:adapted-to-their-world:1}
Uma característica do corpo ou do comportamento que encaixa um \omterm{def:g1:living-or-not:living}{ser
vivo} nas condições do \omterm{def:g2:caring-for-nature:environment}{ambiente} dele. \omterm{def:g1:animals-around-us:animal}{Animal}: o isolamento duplo de
pelo e gordura do urso-polar. \omterm{def:g1:plants-around-us:plant}{Planta}: o \omterm{def:g1:plants-around-us:parts}{caule} do cacto que armazena
água.
\end{solution}

\begin{solution}{exo:g4:adapted-to-their-world:2}
Falta de água: beber cem litros de uma vez e durar dias sem beber,
além da reserva de gordura das corcovas. Areia fofa: pés largos de
dois dedos. Poeira levantada pelo vento: narinas que se fecham e
cílios compridos e duplos. (Calor e fome: de novo as corcovas ---
gordura para travessias magras.)
\end{solution}

\begin{solution}{exo:g4:adapted-to-their-world:3}
Frio: pelo grosso sobre uma camada funda de gordura. Água gelada e
neve: patas largas e um pouco palmadas --- raquetes de neve e remos.
Presas desconfiadas sobre o gelo branco: a camuflagem de uma pelagem
branca.
\end{solution}

\begin{solution}{exo:g4:adapted-to-their-world:4}
Os espinhos são \omterm{def:g1:plants-around-us:parts}{folhas} encolhidas: eles quase não perdem água e
defendem a \omterm{def:g1:plants-around-us:plant}{planta} contra os \omterm{def:g1:animals-around-us:animal}{animais} sedentos. O \omterm{def:g1:plants-around-us:parts}{caule} em barril
armazena água e assumiu o trabalho de captar luz que era das \omterm{def:g1:plants-around-us:parts}{folhas}.
\end{solution}

\begin{solution}{exo:g4:adapted-to-their-world:5}
A forma de almofada quebra o vento gelado: dentro dela o ar fica mais
ameno. A própria forma é o abrigo.
\end{solution}

\begin{solution}{exo:g4:adapted-to-their-world:6}
Caçados: cores de camuflagem, armadura ou espinhos, velocidade,
\omterm{def:g1:the-five-senses:sense}{olhos} bem afastados vigiando em volta (duas bastam). Caçadores:
\omterm{def:g1:the-five-senses:sense}{olhos} voltados para a frente, sentidos aguçados, garras, velocidade
(duas bastam).
\end{solution}

\begin{solution}{exo:g4:adapted-to-their-world:7}
Pelo denso e macio --- um mamífero; garras dianteiras enormes em
forma de pá --- um cavador; \omterm{def:g1:the-five-senses:sense}{olhos} minúsculos e fracos --- vida no
escuro. Juntando tudo: um mamífero cavador subterrâneo, a toupeira.
\end{solution}

\begin{solution}{exo:g4:adapted-to-their-world:8}
Resposta aberta. Exemplo com o pato: pés palmados (para remar), \omterm{def:g1:animals-around-us:covering}{penas}
que escorrem a água (para ficar seco e quente boiando), bico chato de
peneirar (para se alimentar na água).
\end{solution}

\begin{solution}{exo:g4:adapted-to-their-world:9}
As orelhas são radiadores: as grandes e finas soltam o calor do
corpo. No deserto, soltar calor é o problema a resolver --- orelhas
enormes esfriam a raposa. No gelo, guardar calor é tudo --- orelhas
minúsculas perdem o mínimo. Um órgão, dois climas, dois tamanhos
opostos.
\end{solution}

\begin{solution}{exo:g4:adapted-to-their-world:10}
Corpo hidrodinâmico, pés palmados e fôlego dizem: um nadador e
mergulhador; o pelo denso que escorre a água diz: um mamífero que
vive metade da vida em água fria; os bigodes dizem: tateando o
caminho em água escura ou turva. Um mamífero de rios e lagos, que
vive na água e constrói represas --- o castor.
\end{solution}

\begin{solution}{exo:g4:adapted-to-their-world:11}
Foi a água. A baleia e a truta não herdaram nada com formato de
nadador de uma família em comum --- uma é mamífera, a outra é peixe
--- e, mesmo assim, o mesmo problema, mover-se depressa dentro da
água, favorece a mesma resposta hidrodinâmica nas duas. As \omterm{def:g4:adapted-to-their-world:adaptation}{adaptações}
respondem aos problemas do \omterm{def:g2:caring-for-nature:environment}{ambiente}; problemas parecidos podem
extrair formas parecidas de famílias muito diferentes.
\end{solution}
